\documentclass[11pt]{combiArticle}
\usepackage{tikz,amsmath,amsfonts,latexsym}
\usepackage{amssymb}
\usetikzlibrary{shapes,arrows,calc}
\usepackage{calc}
\usepackage{parskip}

\definecolor{myblue}{RGB}{0, 84, 159}

\lang{de}

% Abkuerzungen fuer die Notation
\newcommand{\Dm}{D^{-}}
\newcommand{\Db}{\bar{D}}
\newcommand{\ind}[1]{\mathbf{1}\!\left[#1\right]}
\newcommand{\code}[1]{\texttt{#1}}

\begin{document}
\course{Schichtplan Manager}
\assistent{Linus Palm}
\contact{palm@combi.rwth-aachen.de}
\title{Modell}
\date{04. November 2026}
\maketitle

Dieses Dokument beschreibt das gemischt-ganzzahlige lineare Programm (MILP), mit dem der Gurobi-Solver
(\code{api/app/solver/milp/service.py}, Klasse \code{GurobiSolverService}) den Dienstplan eines Monats
erstellt. Jeder Abschnitt gibt zuerst die mathematische Formulierung und darunter die Erklärung mit Verweis
auf die Gleichungen. In Klammern steht jeweils der Name im Code. Werte, die im Projekt fest gesetzt sind,
stehen bei den Parametern; Werte, die aus TimeOffice bzw. der Webapp kommen, sind Eingabedaten des Monats.

%==========================================================================================================
\section{Mengen und Indizes}

\begin{align}
  & i \in I && \text{Mitarbeitende} \label{set:I}\\
  & u \in U,\quad U^{St} \subseteq U && \text{Planungseinheiten, davon Stationen} \label{set:U}\\
  & s \in S,\quad S^{F}, S^{S}, S^{N} \subseteq S && \text{Schichten, davon Früh-, Spät- und Nachtdienste} \label{set:S}\\
  & l \in L = \{\text{Fachkraft}, \text{Hilfskraft}, \text{Azubi}, \text{MFA}\} && \text{Qualifikationen} \label{set:L}\\
  & d \in D = \{1,\dots,n\} && \text{Tage des Planungsmonats} \label{set:D}\\
  & d \in \Dm = \{1-m,\dots,0\},\quad \Db = \Dm \cup D && \text{Vortage mit bekannten Diensten, Planungshorizont} \label{set:Dm}\\
  & D^{Sa} \subseteq \{0,\dots,n-1\} && \text{Samstage, deren Sonntag im Monat liegt} \label{set:DSa}\\
  & W = (w_1,\dots,w_{|W|}) \subseteq \{-m,\dots,n\} && \text{Samstage der Wochenenden mit einem Tag in } \Db \label{set:W}\\
  & D^{end} = \{1-\max(\Delta^N, L^N+1),\dots,n-1\} && \text{Tage, an denen ein Nachtblock enden kann} \label{set:Dend}\\
  & K \subseteq S \times D \times U^{St} \times L && \text{Bedarfszeilen mit } r_{sdul} > 0 \label{set:K}\\
  & A \subseteq I \times S \times D \times U^{St} && \text{zulässige Dienste, siehe \eqref{eq:A}} \label{set:A}
\end{align}

\textbf{Erklärung.}
\eqref{set:I}--\eqref{set:L}: $I$ (\code{I}) sind alle Mitarbeitenden der gewählten Stationen und ihrer
Springerpools, $U$ (\code{U}) alle Planungseinheiten. Dienste gibt es nur auf Stationen $U^{St}$
(\code{dataset.station\_ids}); Springerpools sind nur Herkunft. Die Schichtmengen nach Typ
(\code{S\_early}, \code{S\_late}, \code{S\_night}) sorgen dafür, dass jede Regel für alle Schichten eines
Typs gilt, unabhängig vom Kürzel.
\eqref{set:D}, \eqref{set:Dm}: Tage werden fortlaufend nummeriert. $D$ (\code{D}) ist der Monat, $\Dm$
(\code{D\_prev}) sind die Vortage, für die die tatsächlichen Dienste bekannt sind (in der Regel $m = 5$). Auf
diesen Tagen gibt es keine Entscheidungsvariablen, die Dienste dort sind feste Werte, siehe \eqref{eq:X}.
\eqref{set:DSa}, \eqref{set:W}: $D^{Sa}$ (\code{D\_s}) enthält die Samstage, deren Sonntag im Monat liegt;
$W$ (\code{W}) jedes Wochenende, dessen Samstag oder Sonntag im Monat oder in den Vortagen liegt, durch seinen
Samstag und aufsteigend; damit gehört das letzte Wochenende vor dem Monat dazu, gleich an welchem Wochentag der
Monat beginnt.
Ist der letzte Vortag ein Samstag, gehört er für die Personen zu $D^{Sa}$, die an ihm gearbeitet haben.
\eqref{set:Dend}: $D^{end}$ (\code{D\_end}) sind die Tage, für die ein Blockende in den Nachtregeln
gebraucht wird.
Allgemein gilt: Alle Tagesbereiche, die in die Vortage reichen (nicht-positive $d$), sind mit $\Dm$
geschnitten, umfassen also nur bekannte Vortage.
\eqref{set:K}: $K$ (Schlüssel von \code{R\_min}) enthält nur Zeilen mit tatsächlichem Bedarf.

%==========================================================================================================
\section{Parameter}

\subsection{Eingabedaten des Monats}

\begin{align}
  & r_{sdul} \in \mathbb{N} && \text{Mindestbesetzung von Schicht } s \text{ am Tag } d \text{ auf Station } u \text{ mit Qualifikation } l \label{par:r}\\
  & h_s \in \mathbb{N} && \text{bezahlte Minuten der Schicht } s \label{par:h}\\
  & t_i,\ c_i \in \mathbb{N} && \text{Sollminuten und gutgeschriebene Minuten im Monat} \label{par:t}\\
  & \lambda_{iud} \in L && \text{Qualifikation der an Tag } d \text{ gültigen Zugehörigkeit von } i \text{ zu } u \label{par:lambda}\\
  & y^{0}_{isd} \in \{0,1\},\ d \in \Dm && \text{fester Dienst } s \text{ von } i \text{ an Vortag } d \label{par:y0}\\
  & T_i \subseteq D && \text{Fortbildungstage von } i \label{par:T}\\
  & B \subseteq I \times D^{+},\ V \subseteq B && \text{blockierte Tage, davon Urlaubstage} \label{par:B}\\
  & \operatorname{home}(i,d) \in U && \text{Stammeinheit von } i \text{ am Tag } d \label{par:home}
\end{align}

Abgeleitet aus dem Bedarf:
\begin{align}
  & R_{usd} = \textstyle\sum_{l} r_{sdul},\quad R_{ud} = \sum_{s} R_{usd},\quad R_{u} = \sum_{d \in D} R_{ud} \label{par:R}\\
  & U^{D} = \{u \in U^{St} : R_u > 0\} \label{par:UD}\\
  & \mathcal{B} = \{(u,s,l) : \exists\, d \in D \text{ mit } r_{sdul} > 0\} \label{par:demanded}\\
  & I_{udl} = \{i \in I : \lambda_{iud} = l\} \label{par:Iudl}
\end{align}

\textbf{Erklärung.}
\eqref{par:r} (\code{R\_min}) stammt aus der Seite „Mindestbesetzung“.
\eqref{par:h} (\code{h}) sind die Netto-Minuten der Schicht, die auf das Monatskonto gebucht werden. In den
Beispieldaten: F 05:55--13:25 (420\,min), Z 08:30--14:15 (345\,min), S 13:15--21:00 (435\,min),
N 20:10--06:10 (555\,min).
\eqref{par:t} (\code{t}, \code{credits}): Soll und Gutschriften kommen aus TimeOffice; für Urlaub und
Fortbildung, die in der Webapp eingetragen werden, wird vorübergehend das Tagessoll gutgeschrieben.
\eqref{par:lambda} (\code{level\_at}): Ein Dienst zählt mit der Qualifikation der Stationszugehörigkeit,
nicht mit der der Person.
\eqref{par:y0} (\code{past\_shifts}) sind die Dienste der Vortage aus dem Kontext.
\eqref{par:T} (\code{training}), \eqref{par:B} (\code{blocked}, \code{vacation}): $B$ enthält alle Tage
mit Urlaub, freiem Tag, Fortbildung oder „nicht verfügbar“, $V$ nur die Urlaubstage. Beide umfassen auch den
Tag $n+1$ nach dem Monat (\code{context.availability}), in den eine Nacht vom letzten Monatstag hineinreicht:
$D^{+} = D \cup \{n+1\}$.
\eqref{par:home} (\code{home\_of}) ist die Einheit der an $d$ gültigen Stammzugehörigkeit.
\eqref{par:R}--\eqref{par:Iudl}: $R_{usd}$, $R_{ud}$, $R_u$ (\code{required\_usd}, \code{required\_ud},
\code{required\_u}) sind der Bedarf je Schicht, Tag und Monat über alle Qualifikationen. $\mathcal{B}$
(\code{demanded}) enthält die Kombinationen aus Station, Schicht und Qualifikation, die im Monat überhaupt
Bedarf haben. $I_{udl}$ (\code{staff\_for}) sind die Personen, die an $d$ auf $u$ als $l$ zählen.

\subsection{Im Projekt gesetzte Parameter}

\begin{center}
\small
\begin{tabular}{l l l p{5.4cm}}
\hline
Symbol & Code & Wert & Bedeutung \\
\hline
$\tau$ & \code{POLICY.balance\_tolerance\_minutes} & 460 & erlaubte Saldoabweichung in Minuten (7:40\,h) \\
$K^{S}$ & \code{K\_s} & 10 & höchstens so viele Arbeitstage in Folge \\
$N^{max}$ & \code{POLICY.max\_consecutive\_nights} & 3 & höchstens so viele Nächte in Folge \\
$L^{N}$ & \code{L\_n} & $\lfloor 2880/1440 \rfloor = 2$ & freie Tage nach einem Nachtblock \\
$\Delta^{N}$ & \code{Dist\_n} & 7 & Tage zwischen Ende eines Nachtblocks und Beginn des nächsten \\
$\pi$ & \code{transfer\_penalty} & 10 & Strafe für einen Dienst auf fremder Station \\
$\omega$ & \code{wish\_reward} & 30 & Gewicht eines Wunsches \\
\hline
\end{tabular}
\end{center}

\subsection{Präferenzwerte}

Für jeden zulässigen Dienst $(i,s,d,u) \in A$ (\code{pref}):
\begin{align}
  p_{isdu} = {} & - \pi \cdot \ind{\operatorname{home}(i,d) \in U^{St},\ \operatorname{home}(i,d) \neq u}
                 \label{par:p}\\
  & - \omega \cdot \ind{\text{Wunsch „freier Tag“ für } (i,d)}
    - \omega \cdot \ind{s \in S^{N},\ \text{Wunsch „freier Tag“ für } (i,d+1)} \notag\\
  & - \omega \cdot \ind{\text{Wunsch „Schicht } s \text{ frei“ für } (i,d)} \notag\\
  & + \omega \cdot \ind{\text{Wunsch „Wunschtag“ für } (i,d)}
    + \omega \cdot \ind{\text{Wunsch „Wunschschicht } s\text{“ für } (i,d)} \notag
\end{align}

\textbf{Erklärung.}
\eqref{par:p}: Ein Dienst auf einer fremden Station kostet $\pi$, aber nur, wenn die Person eine Station als
Stammeinheit hat; Springer (Stammeinheit ist ein Pool) dürfen überall gleich gut arbeiten. Ein erfüllter
Dienstwunsch bringt $+\omega$, ein Dienst gegen einen Freiwunsch kostet $-\omega$. Beim freien Tag zählt auch
ein Nachtdienst am Vorabend, weil er bis 06:10 in den Wunschtag hineinreicht; so bewertet es auch der
Schedule-Check.

%==========================================================================================================
\section{Zulässige Dienste}

\begin{align}
  N^{V} &= \{(i,s,d) : \text{für } (i,d) \text{ ist eine Einschränkung eingetragen, die } s \text{ nicht erlaubt}\} \label{eq:NV}\\
  N^{N} &= \{(i,s,d) : s \in S^{N},\ d \in D,\ (i,d+1) \in B\} \label{eq:NN}\\
  N^{S} &= \{(i,s,d) : s \in S^{S},\ d \in D,\ (i,d+1) \in V,\ (i,d) \notin V\} \label{eq:NS}\\
  A &= \bigl\{(i,s,d,u) : u \in U^{St},\ \lambda_{iud} \text{ definiert},\
       (i,s,d) \notin N^{V} \cup N^{N} \cup N^{S},\ (u,s,\lambda_{iud}) \in \mathcal{B}\bigr\} \label{eq:A}
\end{align}

\textbf{Erklärung.}
Statt diese Regeln als Nebenbedingungen zu formulieren, werden verbotene Dienste gar nicht erst als Variable
angelegt (\code{not\_allowed}, Liste beim Anlegen von \code{assignment}).
\eqref{eq:NV}: Urlaub, freier Tag, Fortbildung und „nicht verfügbar“ verbieten alle Schichten, „nur
bestimmte Schichten“ verbietet die übrigen.
\eqref{eq:NN}: Eine Nacht reicht in den Folgetag; vor einem blockierten Tag ist sie verboten, am letzten
Monatstag also auch vor einer Abwesenheit am Ersten des Folgemonats.
\eqref{eq:NS}: Am Tag vor Urlaubsbeginn gibt es keinen Spätdienst, auch wenn der Urlaub am Ersten des
Folgemonats beginnt.
\eqref{eq:A}: Ein Dienst ist nur auf einer Station möglich, zu der an diesem Tag eine gültige Zugehörigkeit
besteht, und nur in einer Schicht, für die die Station mit dieser Qualifikation im Monat überhaupt Bedarf hat
(z.\,B. kein Zwischendienst, wenn nie einer gefordert ist, keine Nächte für Azubis, wenn nie Azubi-Nachtbedarf
besteht).

%==========================================================================================================
\section{Entscheidungsvariablen}

\begin{align}
  & x_{isdu} \in \{0,1\} && (i,s,d,u) \in A && \text{Dienst } (\code{assignment}) \label{var:x}\\
  & g_{sdul} \ge 0 && (s,d,u,l) \in K && \text{Lücke } (\code{gap}) \label{var:g}\\
  & \eta_i \in \mathbb{R} && i \in I && \text{Monatssaldo in Minuten } (\code{hour}) \label{var:eta}\\
  & o_i \ge 0 && i \in I && \text{Saldo außerhalb des Bands } (\code{over}) \label{var:o}\\
  & q_{ud} \ge 0 && u \in U^{St},\ d \in D,\ R_{ud} > 0 && \text{Quadrat der Lücken des Tages } (\code{gap\_sq}) \label{var:q}\\
  & \rho_{ud} \ge 0 && u \in U^{D},\ d \in D && \text{Abweichung der Tagesbesetzung } (\code{day\_dev}) \label{var:rho}\\
  & \delta_{usd} \ge 0 && u \in U^{St},\ d \in D,\ R_{ud} > 0,\ s \in S && \text{Abweichung der Schichtbesetzung } (\code{dev}) \label{var:delta}\\
  & b_{id} \in [0,1] && i \in I,\ d \in \{0\} \cup D && \text{Beginn eines Nachtblocks } (\code{n\_beg}) \label{var:b}\\
  & e_{id} \in [0,1] && i \in I,\ d \in D^{end} && \text{Ende eines Nachtblocks } (\code{n\_end}) \label{var:e}\\
  & z_{id} \in [0,1] && i \in I,\ d \in D && \text{Spät- vor Frühdienst } (\code{L\_to\_E}) \label{var:z}\\
  & v_{isd} \in [0,1] && i \in I,\ s \in S,\ d \in D && \text{gleiche Schicht wie am Vortag } (\code{cons\_s}) \label{var:v}
\end{align}

\textbf{Erklärung.}
Nur $x$ \eqref{var:x} ist ganzzahlig. Alle anderen Variablen sind stetige Hilfsgrößen; ihre Werte ergeben
sich aus $x$ und der Optimierung. $g$ \eqref{var:g} nimmt im Optimum ganzzahlige Werte an, $b$, $e$, $z$
und $v$ \eqref{var:b}--\eqref{var:v} nur 0 oder 1, weil sie ausschließlich durch ganzzahlige Ausdrücke in $x$
nach oben bzw. unten begrenzt werden und die Optimierung sie an diese Grenze drückt.

%==========================================================================================================
\section{Hilfsausdrücke}

\begin{align}
  X_{isd} &= \begin{cases}
      \sum_{u : (i,s,d,u) \in A} x_{isdu} & d \in D\\
      y^{0}_{isd} & d \in \Dm\\
      0 & \text{sonst}
    \end{cases} && \text{(\code{x}, \code{shift\_on})} \label{eq:X}\\
  Y_{id} &= \textstyle\sum_{s \in S} X_{isd} && \text{(\code{work})} \label{eq:Y}\\
  Y^{F}_{id} &= \textstyle\sum_{s \in S^{F}} X_{isd},\quad
  Y^{S}_{id} = \sum_{s \in S^{S}} X_{isd},\quad
  Y^{N}_{id} = \sum_{s \in S^{N}} X_{isd} && \text{(\code{early}, \code{late}, \code{night})} \notag\\
  P_{usd} &= \textstyle\sum_{i : (i,s,d,u) \in A} x_{isdu} && \text{(\code{staffed})} \label{eq:P}\\
  P_{ud} &= \textstyle\sum_{s} P_{usd},\quad
  P_{u} = \sum_{d \in D} P_{ud} && \text{(\code{staff\_of\_day}, \code{staff\_of\_month})} \notag\\
  G_{ud} &= \textstyle\sum_{(s,d,u,l) \in K} g_{sdul} && \text{(\code{gaps\_of\_day})} \label{eq:G}
\end{align}

\textbf{Erklärung.}
\eqref{eq:X}: $X_{isd}$ ist 1, wenn $i$ an $d$ Schicht $s$ arbeitet, egal auf welcher Station. Auf den
Vortagen ist es der feste Dienst aus dem Kontext, außerhalb des bekannten Zeitraums 0. Dadurch wirken alle
Regeln über den Monatswechsel, ohne dass die Vergangenheit verändert werden kann.
\eqref{eq:Y}: $Y_{id}$ ist 1, wenn $i$ an $d$ irgendeinen Dienst hat, $Y^{F}$, $Y^{S}$, $Y^{N}$ entsprechend
für Früh-, Spät- und Nachtdienst.
\eqref{eq:P}: Besetzung einer Station je Schicht, je Tag und im Monat.
\eqref{eq:G}: Summe der Lücken einer Station an einem Tag.

%==========================================================================================================
\section{Zielfunktionen}

Die Ziele werden lexikografisch in dieser Reihenfolge optimiert (Priorität fallend):
\begin{align}
  f_1 &= \textstyle\sum_{(s,d,u,l) \in K} g_{sdul} \to \min && \text{(\code{gaps})} \label{obj:gaps}\\
  f_2 &= \textstyle\sum_{i \in I} o_i \to \min && \text{(\code{balance\_violation})} \label{obj:balance}\\
  f_3 &= \textstyle\sum_{(u,d) : R_{ud} > 0} \dfrac{q_{ud}}{R_{ud}} \to \min && \text{(\code{gap\_spread})} \label{obj:gapspread}\\
  f_4 &= \textstyle\sum_{u \in U^{D}} \sum_{d \in D} \rho_{ud} \to \min && \text{(\code{staffing\_spread})} \label{obj:staffingspread}\\
  f_5 &= \textstyle\sum_{(i,s,d,u) \in A} p_{isdu}\, x_{isdu} \to \max && \text{(\code{wishes})} \label{obj:wishes}\\
  f_6 &= \textstyle\sum_{u,s,d} \delta_{usd} \to \min && \text{(\code{shift\_balance})} \label{obj:shiftbalance}\\
  f_7 &= \textstyle\sum_{i \in I} \sum_{d \in D} z_{id} \to \min && \text{(\code{Min\_L\_to\_E})} \label{obj:ltoe}\\
  f_8 &= \textstyle\sum_{i \in I} \sum_{s \in S} \sum_{d \in D} v_{isd} \to \max && \text{(\code{cons})} \label{obj:cons}
\end{align}

\textbf{Lexikografisches Vorgehen.}
Für $k = 1,\dots,8$ wird $f_k$ über alle zulässigen Lösungen optimiert, unter den zusätzlichen Bedingungen
\begin{equation}
  f_j \le \hat f_j + \alpha_j \quad (\text{bzw. } f_j \ge \hat f_j - \alpha_j \text{ für } j = 5, 8)
  \qquad \forall j < k, \label{obj:lex}
\end{equation}
wobei $\hat f_j$ der in Stufe $j$ erreichte Wert ist. Die Toleranzen $\alpha_j$ (\code{abstol}) sind
\begin{center}
\small
\begin{tabular}{l l l}
\hline
Stufe & Gesamtspielraum $\varepsilon_j$ (\code{stage\_slack}) & $\alpha_j = \varepsilon_j/2$ \\
\hline
$f_3$ \code{gap\_spread} & $2$ & $1$ \\
$f_4$ \code{staffing\_spread} & $|U^{D}| \cdot n / 7$ & $|U^{D}| \cdot n / 14$ \\
$f_6$ \code{shift\_balance} & $2\,|U^{D}| \cdot n / 7$ & $|U^{D}| \cdot n / 7$ \\
alle anderen & -- & $10^{-6}$ \\
\hline
\end{tabular}
\end{center}

\textbf{Abbruchkriterien.} Gesamtzeit $T$ (in der Demo-Version fest $T = 30$\,s, \code{SOLVER\_MAX\_TIME\_SECONDS}),
relative Optimalitätslücke 3\,\% (\code{MIPGap}) für jede Stufe außer $f_1$ und $f_2$, die exakt gelöst
werden (\code{MIPGap} 0); für $f_3$, $f_4$, $f_6$ zusätzlich eine absolute
Lücke von $\varepsilon_j/2$ (\code{MIPGapAbs}); Zeitlimit $0{,}15\,T$ für $f_4$ und $0{,}10\,T$ für $f_6$
(\code{stage\_time\_share}). Ein Watchdog beendet die Rechnung spätestens 5\,s nach $T$, gleich in welcher Stufe.

\textbf{Erklärung.}
\eqref{obj:gaps}: Zuerst möglichst wenige unbesetzte Pflichtstellen.
\eqref{obj:balance}: Danach die Saldoregel: Minuten, um die ein Monatssaldo das Band $\pm\tau$ verlässt,
siehe \eqref{con:over}. Die Regel ist weich, damit es auch dann einen Plan gibt, wenn sie nicht für alle
einzuhalten ist.
\eqref{obj:gapspread}: Lücken möglichst gleichmäßig auf die Tage verteilen. $q_{ud}$ ist das Quadrat der
Lücken eines Tages \eqref{con:gapsq}; zwei Tage mit je einer Lücke ($1+1$) sind billiger als ein Tag mit zwei
($4$). Die Division durch $R_{ud}$ gewichtet relativ zum Bedarf: Eine Lücke wiegt an einem schwach besetzten
Tag schwerer.
\eqref{obj:staffingspread}: Jeder Tag einer Station soll im gleichen Verhältnis Ist/Soll besetzt sein wie der
ganze Monat \eqref{con:daydev}; Über- und Unterbesetzung verteilen sich so gleichmäßig auf die Tage.
\eqref{obj:wishes}: Wünsche erfüllen und Einsätze auf fremden Stationen vermeiden, siehe \eqref{par:p}.
\eqref{obj:shiftbalance}: Die Besetzung eines Tages soll sich im Verhältnis der Mindestbesetzung auf die
Schichten verteilen \eqref{con:shiftdev}.
\eqref{obj:ltoe}: Möglichst wenige Wechsel von Spät- auf Frühdienst \eqref{con:ltoe}; die 11-h-Ruhezeit ist
im Modell bewusst nur weich.
\eqref{obj:cons}: Möglichst oft dieselbe Schicht wie am Vortag \eqref{con:cons}.
\eqref{obj:lex}: Eine spätere Stufe darf eine frühere nur um deren Toleranz verschlechtern. Bei den drei
Verteilungsstufen ist das bewusst etwas Spielraum, damit die späteren Stufen noch Wahlmöglichkeiten haben;
die andere Hälfte des Spielraums erlaubt einen frühen Abbruch der Stufe. Insgesamt liegt der Endwert einer
solchen Stufe höchstens $\varepsilon_j$ über dem bestmöglichen Wert.

%==========================================================================================================
\section{Nebenbedingungen}

\subsection{Grundregeln und Besetzung}

\begin{align}
  & Y_{id} \le 1 && \forall i \in I,\ d \in D \label{con:one}\\
  & \textstyle\sum_{i \in I_{udl} : (i,s,d,u) \in A} x_{isdu} + g_{sdul} \ge r_{sdul} && \forall (s,d,u,l) \in K \label{con:demand}
\end{align}

\textbf{Erklärung.}
\eqref{con:one} (\code{one\_shift}): Höchstens ein Dienst pro Person und Tag, über alle Stationen.
\eqref{con:demand} (\code{demand}): Die Mindestbesetzung jeder Bedarfszeile wird durch Dienste mit passender
Qualifikation gedeckt; was fehlt, ist eine Lücke $g$. Die Mindestbesetzung ist damit weich und wird über
\eqref{obj:gaps} minimiert.

\subsection{Monatssaldo}

\begin{align}
  & \textstyle\sum_{(s,d,u) : (i,s,d,u) \in A} h_s\, x_{isdu} + c_i - \eta_i = t_i && \forall i \in I \label{con:balance}\\
  & o_i \ge \eta_i - \tau,\qquad o_i \ge -\eta_i - \tau && \forall i \in I \label{con:over}
\end{align}

\textbf{Erklärung.}
\eqref{con:balance} (\code{workhours}): $\eta_i$ ist der Saldo in Minuten: geplante Minuten plus
Gutschriften minus Soll, genau wie das Monatskonto im Projekt.
\eqref{con:over} (\code{over\_up}, \code{over\_down}): $o_i$ misst, um wie viele Minuten der Saldo außerhalb
von $\pm\tau = \pm 460$ liegt. Im Optimum gilt $o_i = \max(0, |\eta_i| - \tau)$.

\subsection{Gleichmäßige Verteilung}

\begin{align}
  & q_{ud} \ge (2k+1)\, G_{ud} - k(k+1) && \forall (u,d) : R_{ud} > 0,\ k = 0,\dots,R_{ud}-1 \label{con:gapsq}\\
  & \rho_{ud} \ge P_{ud} - \tfrac{R_{ud}}{R_u}\, P_u,\qquad
    \rho_{ud} \ge -\bigl(P_{ud} - \tfrac{R_{ud}}{R_u}\, P_u\bigr) && \forall u \in U^{D},\ d \in D \label{con:daydev}\\
  & \delta_{usd} \ge P_{usd} - \tfrac{R_{usd}}{R_{ud}}\, P_{ud},\qquad
    \delta_{usd} \ge -\bigl(P_{usd} - \tfrac{R_{usd}}{R_{ud}}\, P_{ud}\bigr) && \forall (u,d) : R_{ud} > 0,\ s \in S \label{con:shiftdev}
\end{align}

\textbf{Erklärung.}
\eqref{con:gapsq} (\code{gap\_sq}): Lineare Darstellung von $q_{ud} \ge G_{ud}^2$. Die rechte Seite ist für
jedes $k$ die Sekante durch die Parabelpunkte $(k, k^2)$ und $(k+1, (k+1)^2)$. Zwischen $k$ und $k+1$ liegt sie
über der Parabel, außerhalb darunter. Das Maximum aller Sekanten ist daher die stückweise lineare Interpolation
von $G_{ud}^2$: überall mindestens $G_{ud}^2$ und in ganzzahligen Punkten genau $G_{ud}^2$. Da die Lücken eines
Tages ganzzahlig zwischen $0$ und $R_{ud}$ liegen und $q_{ud}$ minimiert wird, ist $q_{ud} = G_{ud}^2$ im
Optimum exakt.
\eqref{con:daydev} (\code{day\_dev\_up}, \code{day\_dev\_down}): $\rho_{ud} \ge |P_{ud} - \frac{R_{ud}}{R_u}
P_u|$, also die Abweichung der Tagesbesetzung von ihrem Anteil an der Monatsbesetzung. Ein Tag ohne Bedarf
hat Anteil 0, jeder Dienst dort zählt als Abweichung.
\eqref{con:shiftdev} (\code{shift\_dev\_up}, \code{shift\_dev\_down}): $\delta_{usd} \ge |P_{usd} -
\frac{R_{usd}}{R_{ud}} P_{ud}|$, die Abweichung einer Schicht von ihrem Anteil an der Tagesbesetzung. Weil
sich Anteile nur mit „Bruchteilen von Personen“ exakt treffen lassen, ist ein Teil dieser Abweichung
unvermeidbar; die Stufe hat deshalb ein Zeitlimit statt eines Optimalitätsbeweises.

\subsection{Folgedienste und Wechsel}

\begin{align}
  & v_{isd} \le X_{isd},\qquad v_{isd} \le X_{is,d-1} && \forall i \in I,\ s \in S,\ d \in D \label{con:cons}\\
  & z_{id} \ge Y^{S}_{i,d-1} + Y^{F}_{id} - 1 && \forall i \in I,\ d \in D \label{con:ltoe}\\
  & \textstyle\sum_{d' = d}^{d + K^{S}} Y_{id'} + |T_i \cap \{d,\dots,d+K^{S}\}| \le K^{S}
    && \forall i \in I,\ d \in \{1-\min(m,K^{S}),\dots,0\} \cup \{1,\dots,n-K^{S}\} \label{con:ks}
\end{align}

\textbf{Erklärung.}
\eqref{con:cons} (\code{cons\_s}): $v_{isd}$ darf nur 1 sein, wenn an $d$ und $d-1$ dieselbe Schicht gearbeitet
wird; maximiert in \eqref{obj:cons}.
\eqref{con:ltoe} (\code{L\_to\_E}): $z_{id}$ ist mindestens 1, wenn auf einen Spätdienst ein Frühdienst
folgt; minimiert in \eqref{obj:ltoe}.
\eqref{con:ks}: In jedem Fenster von $K^{S}+1 = 11$ Tagen höchstens $K^{S} = 10$ Arbeitstage; Fortbildungstage
zählen als Arbeitstage. Die Fenster beginnen schon in den Vortagen, damit die Regel über den Monatswechsel gilt.

\subsection{Nachtdienste}

\begin{align}
  & \textstyle\sum_{d' = d}^{d + N^{max}} Y^{N}_{id'} \le N^{max}
    && \forall i \in I,\ d \in \{1-\min(m,N^{max}),\dots,0\} \cup \{1,\dots,n-N^{max}\} \label{con:nights}\\
  & b_{id} \ge Y^{N}_{id} - Y^{N}_{i,d-1} && \forall i \in I,\ d \in \{0\} \cup D \label{con:beg}\\
  & e_{id} \ge Y^{N}_{id} - Y^{N}_{i,d+1} && \forall i \in I,\ d \in D^{end} \label{con:end}\\
  & b_{id} + e_{id} \le 1 && \forall i \in I,\ d \in \{0,\dots,n-1\} \label{con:single}\\
  & e_{id} + Y_{i,d+j} \le 1 && \forall i \in I,\ d \in \{-L^{N},\dots,n-1\},\ j = 1,\dots,L^{N},\ d+j \in D \label{con:recovery}\\
  & e_{id} + Y^{F}_{i,d+L^{N}+1} \le 1 && \forall i \in I,\ d \in \{-L^{N},\dots,n-1\},\ d+L^{N}+1 \in D \label{con:recoveryearly}\\
  & e_{id} + \textstyle\sum_{j = 1,\ d+j \in D}^{\Delta^{N}} b_{i,d+j} \le 1 && \forall i \in I,\ d \in \{1-\Delta^{N},\dots,n-1\} \cap D^{end} \label{con:distance}
\end{align}

\textbf{Erklärung.}
\eqref{con:nights}: In jedem Fenster von $N^{max}+1 = 4$ Tagen höchstens $N^{max} = 3$ Nächte.
\eqref{con:beg}, \eqref{con:end} (\code{n\_beg}, \code{n\_end}): $b_{id}$ ist mindestens 1, wenn an $d$ ein
Nachtblock beginnt (Nacht an $d$, keine an $d-1$), $e_{id}$, wenn einer endet. Untergrenzen genügen, weil $b$
und $e$ nur in $\le 1$-Bedingungen mit positivem Vorzeichen und in keinem Ziel vorkommen.
\eqref{con:single}: Beginn und Ende am selben Tag hieße Ein-Nacht-Block; der ist verboten. Für den letzten
Vortag erzwingt das eine zweite Nacht am Monatsersten, wenn der Vortag eine einzelne Nacht war.
\eqref{con:recovery}, \eqref{con:recoveryearly}: Nach dem letzten Nachtdienst eines Blocks $L^{N} = 2$ Tage
frei und am dritten Tag kein Frühdienst. Das bildet die 48 Stunden Erholung ab: Die Nacht endet um 06:10, der
Frühdienst am dritten Tag beginnt um 05:55. Nur Tage im Monat werden eingeschränkt, die Vergangenheit nie
gegen sich selbst.
\eqref{con:distance}: Nach dem Ende eines Nachtblocks darf in den nächsten $\Delta^{N} = 7$ Tagen kein neuer
Block beginnen; berücksichtigt werden nur Blockanfänge im Monat.

\subsection{Wochenenden}

\begin{align}
  & X_{isd} = X_{is,d+1} && \forall i \in I,\ s \in S,\ d \in D^{Sa} \label{con:weekend}\\
  & \Omega_{iw} = \begin{cases}
      \max\bigl(Y_{iw},\, Y_{i,w+1}\bigr) & w+1 \in \Dm\\
      Y_{iw} & w \in \Db,\ w+1 \notin \Dm\\
      Y_{i,w+1} & \text{sonst}
    \end{cases} && \text{(\code{weekend})} \label{eq:Omega}\\
  & \Omega_{i,w_k} + \Omega_{i,w_{k+1}} \le 1 && \forall i \in I,\ k = 1,\dots,|W|-1 \text{ mit } w_{k+1} \ge 0
    \label{con:twoweekends}
\end{align}

\textbf{Erklärung.}
\eqref{con:weekend}: Samstag und Sonntag haben dieselbe Schicht, oder beide sind frei. Ist der letzte Vortag
ein Samstag mit Dienst, erzwingt das denselben Dienst am Monatsersten; war er frei, ist der Sonntag offen. Sonst
müsste in einem Monat, der an einem Sonntag beginnt, jeder am Ersten frei haben, dessen Samstag davor frei war.
\eqref{eq:Omega}: $\Omega_{iw}$ gibt an, ob $i$ das Wochenende mit dem Samstag $w$ arbeitet. Gearbeitet ist es
mit einem Dienst am Samstag oder Sonntag; eine Nacht vom Freitag allein lässt es frei. Im Monat haben Samstag
und Sonntag dieselbe Schicht \eqref{con:weekend}, also entscheidet der Samstag. Ein Wochenende in den Vortagen
ist bekannt, $\Omega$ dort also eine Konstante. Liegt der Samstag vor den bekannten Tagen (Monatsbeginn an einem
Sonntag ohne Vortage), entscheidet der Sonntag.
\eqref{con:twoweekends}: Keine zwei Wochenenden in Folge, auch nicht über den Monatswechsel. Geprüft werden nur
Paare, deren späteres Wochenende den Monat berührt ($w_{k+1} \ge 0$, sein Sonntag liegt also im Monat); die
Vergangenheit wird nie gegen sich selbst geprüft. Daraus
folgt zugleich, dass höchstens jedes zweite Wochenende gearbeitet wird.

%==========================================================================================================
\section{Nach dem Lösen}

Aus der Lösung werden die Dienste $\{(i,s,d,u) : x_{isdu} = 1\}$ mit der Qualifikation $\lambda_{iud}$
übernommen. Die gemeldeten Lücken werden nicht aus $g$ gelesen, sondern aus den Diensten berechnet,
$\max(0,\, r_{sdul} - \sum_{i \in I_{udl}} x_{isdu})$, ebenso die angezeigten Werte von $f_1$, $f_3$, $f_4$ und
$f_6$. So stimmen sie genau mit dem unabhängigen Schedule-Check des Projekts überein, auch wenn Hilfsvariablen
wegen der Toleranzen \eqref{obj:lex} über ihrem echten Wert liegen. Der Schedule-Check prüft den fertigen Plan
anschließend gegen alle harten Regeln des Projekts, darunter auch die 11-h-Ruhezeit, die das Modell nur weich
behandelt \eqref{obj:ltoe}.

\textbf{Zeitumstellung.} Modell und Prüfung ignorieren die Zeitumstellung: Eine Schicht zählt für die tägliche
Höchstarbeitszeit und die Pausen mit ihren geplanten Zeiten, auch die Nacht der Umstellung im Oktober, die
tatsächlich eine Stunde länger dauert.

\textbf{Ersatzruhetag.} Die Demo-Version prüft nicht, ob jeder gearbeitete Sonn- oder Feiertag durch einen freien
Werktag ausgeglichen wird (ArbZG §11 Abs.~3); das Modell bildet diese Regel ebenfalls nicht ab.

\end{document}
